% RESULTADO_RECUPERADO. II REV10,
% 05b_coordenadas_angulares.tex:1--125.
% Se conserva la publicación angular, sus cartas y pruebas; el funcional
% ulterior de Barbero no es una dependencia de la forma elíptica.
\section{Representación angular y orientación de los canales}
\label{iv:angular:section}

La sección \(\hbar_{\mathrm{ret}}\) pertenece a la recta de
acción \(\mathcal L_S\); \(\eta_{\mathrm{ret}}\) es su
coeficiente adimensional en la base declarada en
\eqref{iv:action:hbar}. La representación angular actúa después
sobre \((\alpha,\eta_{\mathrm{ret}})\). Escribiendo
\([t]_{360}\) para la clase de \(t\) en
\(\mathbb R/(360\mathbb Z)\), se define
\begin{equation}
 \mathcal P_{\mathrm{ang}}(\alpha,\eta_{\mathrm{ret}})
 =\bigl([1000\alpha]_{360},
        [2(\eta_{\mathrm{ret}}+\alpha)]_{360}\bigr).
 \label{iv:angular:publicacion}
\end{equation}
La completación nonádica determina el factor \(1000\); la
clausura bidireccional determina el factor dos. La coordenatización circular
de \(360\) unidades conserva las particiones
\(12\cdot30=4\cdot90=3\cdot120\). En la rama positiva, los
representantes levantados son
\begin{equation}
 A_{\deg}=1000\alpha,\qquad
 C^*_{\deg}=2(\eta_{\mathrm{ret}}+\alpha).
 \label{iv:angular:representantes}
\end{equation}
El subíndice distingue estas coordenadas angulares de los bloques
del registro dodecafásico.

\begin{proposition}[Desplazamiento de la representación en base mil]
\label{iv:angular:desplazamiento}
Sea \(\alpha=\sum_{n\geq1}b_n1000^{-n}\), con
\(b_n\in\{0,\ldots,999\}\), la expansión compatible ya
generada. Su coordenada levantada satisface
\[
 A_{\deg}=b_1+\sum_{n\geq1}b_{n+1}1000^{-n}.
\]
La representación conserva los bloques a partir del segundo y
sitúa el bloque inicial en la parte entera.
\end{proposition}
\begin{proof}
La serie converge absolutamente. Multiplicar cada término por
\(1000\) y separar el término inicial proporciona la identidad.
La operación actúa sobre la expansión construida, con sus
bloques y memoria; la representación circular posterior toma
su clase módulo \(360\).
\end{proof}

\subsection{Coordenatizaciones circulares y covariancia de la acción}
\label{iv:angular:cartas}

\begin{definition}[Coordenatizaciones de la coordenada angular]
\label{iv:angular:unidades}
La coordenatización en grados expresa una vuelta completa mediante \(360\)
unidades; la coordenatización en radianes, mediante \(2\pi\). El cambio
de coordenadas y su inversa son
\begin{equation}
 \theta_A=\frac{\pi}{180}A_{\deg},\qquad
 \theta_C=\frac{\pi}{180}C^*_{\deg},\qquad
 A_{\deg}=\frac{180}{\pi}\theta_A,\qquad
 C^*_{\deg}=\frac{180}{\pi}\theta_C.
 \label{iv:angular:cambio-carta}
\end{equation}
Un recorrido con número de vueltas conservado utiliza las
coordenadas levantadas. El cociente circular expresa su
clase módulo la vuelta correspondiente.
\end{definition}

El registro retiene orientación, vueltas y memoria del transporte.
La fase final y el levantamiento son representaciones diferentes:
el segundo distingue recorridos que comparten la primera.
Cada expresión de acción declara cuál de esas coordenadas utiliza.

\begin{proposition}[Covariancia del lector de acción]
\label{iv:angular:accion}
Para una sección \(s\in\{\mathrm{pre},\mathrm{ret}\}\) y una
coordenada angular levantada, se cumple
\begin{equation}
 \mathcal S_s(\theta)
 =\hbar_s\theta_{\mathrm{rad}}
 =h_s\frac{\theta_{\deg}}{360}.
 \label{iv:angular:accion-cartas}
\end{equation}
Una vuelta positiva expresa \(h_s\).
\end{proposition}
\begin{proof}
Sustituir
\(\theta_{\mathrm{rad}}=\pi\theta_{\deg}/180\) y
\(h_s=2\pi\hbar_s\) da
\[
 \hbar_s\frac{\pi\theta_{\deg}}{180}
 =\frac{h_s\theta_{\deg}}{360}.
\]
En una vuelta \(\theta_{\deg}=360\), con lo que resulta \(h_s\).
\end{proof}

La periodicidad de una representación mantiene su condición
propia: si \(U(2\pi)=-I\), la restitución operatoria ocurre
en \(4\pi\). La normalización de la acción por vuelta de
\eqref{iv:angular:accion-cartas} y el retorno del estado
representado conservan así sus dominios respectivos.

\subsection{Canales positivos y reversión de la coordenada impar}
\label{iv:angular:canales}

Los canales se expresan, en cualquiera de las dos coordenatizaciones, mediante
\begin{equation}
 q_\pm=\exp[-(\theta_A\pm\theta_C)]
 =\exp\!\left[-\frac{\pi}{180}
                   (A_{\deg}\pm C^*_{\deg})\right].
 \label{iv:angular:q}
\end{equation}

\begin{proposition}[Recuperación de las coordenadas angulares]
\label{iv:angular:inversion}
En la coordenatización real levantada, los canales positivos determinan
unívocamente
\begin{equation}
 A_{\deg}=-\frac{90}{\pi}\log(q_+q_-),\qquad
 C^*_{\deg}=\frac{90}{\pi}\log\frac{q_-}{q_+}.
 \label{iv:angular:recuperacion}
\end{equation}
El intercambio \((q_+,q_-)\mapsto(q_-,q_+)\) conserva
\(A_{\deg}\) y revierte \(C^*_{\deg}\).
\end{proposition}
\begin{proof}
La suma y la diferencia de los logaritmos de
\eqref{iv:angular:q} son \(-\pi A_{\deg}/90\) y
\(-\pi C^*_{\deg}/90\), respectivamente. El logaritmo real
es unívoco en los canales positivos. Su producto es simétrico
y el cociente se invierte al intercambiarlos.
\end{proof}

La involución antipodal \(r\mapsto r+6\pmod {12}\) del
registro forma seis parejas. La distribución uniforme de la
coordenada impar sobre esas parejas asigna
\(C^*_{\deg}/6\) a cada una; la suma recupera
\(C^*_{\deg}\). La conversión de unidades se aplica una
sola vez a cada argumento. En particular,
\(C^*_{\deg}/A_{\deg}=\theta_C/\theta_A\) cuando
\(A_{\deg}\ne0\), mientras el estado anterior conserva
los ángulos levantados y su memoria.
